% =====================================================================
% Section 1: Introduction — EvidenceFlow
% =====================================================================
%
% Status: Draft §1 Introduction v5 (2026-05-18)
%   v5: advisor-driven 5-paragraph reorganization.
%       Para 1 (~80w):   RAG primer.
%       Para 2 (~200w):  Text RAG → GraphRAG landscape, open issue.
%       Para 3 (~350w):  Failure modes + fig:motivation + framing.
%       Para 4 (~300w):  Method preview (mirrors §4 sub-headings).
%       Para 5 (~150w):  2–3 contribution bullets.
%   v4: see archive/01_intro_v4_pre-advisor-rewrite_20260518.tex
%
% =====================================================================

\section{Introduction}
\label{sec:intro}

Retrieval-augmented generation (RAG) couples a large language model
with a retrieval module that fetches passages from an external
corpus, so that the model can ground its outputs in evidence rather
than rely on parameters alone. RAG is the standard recipe for open
question answering and for any setting where the answer depends on
documents the model was not trained on, because the retrieved
passages provide a check on what the language model says.

The current RAG landscape spans two families of systems. Classic
text RAG ranks each passage independently against the question.
Sparse and dense retrievers such as BM25
\citep{robertson2009bm25} and DPR \citep{karpukhin2020dpr} are the
standard backbones, and hierarchical schemes like RAPTOR
\citep{sarthi2024raptor} cluster and summarize passages into a
retrieval tree to recover broader context. These systems work well
on single-hop questions but tend to miss bridge passages whose
surface similarity to the question is low. A second family,
GraphRAG, organizes the corpus into a structure that exposes
cross-document connections, biasing retrieval toward that structure
rather than toward surface similarity. HippoRAG and HippoRAG~2
\citep{jimenez2024hipporag,gutierrez2025hipporag} run personalized
PageRank over a corpus-wide passage--phrase graph; PropRAG
\citep{wang2025proprag} discovers chains by beam search over a
proposition graph; GraphRAG \citep{edge2024graphrag} builds
community-summary graphs; HGRAG \citep{hgrag2024} represents
entities and passages within a hypergraph. These systems improve
multi-hop coverage at the level of graph relevance, but they still
have to hand the reader a short ranked list of natural-language
passages under a small budget, and graph relevance does not
automatically yield such a list.

\begin{figure}[t]
\centering
\fbox{\parbox{0.92\linewidth}{\centering\vspace{1.5em}
\textit{Two failure modes of converting graph relevance into a top-$K$ reader prefix.}\\[0.3em]
Left panel: bridge underweighting --- a bridge passage on a multi-hop chain receives little PPR mass and falls outside the top-$K$.\\[0.3em]
Right panel: entity-neighborhood crowding --- the top-$K$ list concentrates around one entity neighborhood and leaves other question-side requirements uncovered.
\vspace{1.5em}}}
\caption{Motivating failure modes of GraphRAG retrieval. A useful
top-$K$ list should preserve cross-document bridges and cover the
distinct entities and relations the question asks about. Existing
graph-propagation operators address neither problem on their own.}
\label{fig:motivation}
\end{figure}

These routes all face the same downstream constraint, and two
failure modes recur in this conversion. First, propagation puts
mass on the strongest seed neighborhood, so a top-$K$ list can look
on-topic and still omit the bridge passages that connect to the
answer. Pointwise dense retrieval makes the same mistake from the
opposite direction: it scores passages by question-passage
similarity and therefore underweights bridges whose surface
similarity to the question is low. Second, a list of high-scoring
passages can crowd around a single entity in the question while
leaving other entities or relations entirely uncovered, so the
reader has the ingredients for one hop and nothing for the next.
Path search and summary representations can shift where this
conversion is handled but still leave the retriever with a final
reader-facing ordering decision. Figure~\ref{fig:motivation}
visualizes both failure modes on a stylized multi-hop query.

The problem is therefore not only how to propagate relevance over a
graph, but how to turn propagated relevance into evidence that is
both \emph{connected}---preserving the bridge relations that span
documents---and \emph{covering}, in the sense that the selected
passages collectively address the distinct question-side entities
and relations a reader has to resolve. We treat this conversion as
the core retrieval problem and keep the interface simple. The
system outputs a ranked evidence list $R_q$ of natural-language
passages, not a symbolic path or a generated reasoning chain, and
the reader consumes a prefix of $R_q$.

We propose \methodname{}, a graph-based retrieval framework
operationalized through a staged retrieval pipeline. Offline
indexing builds a fact-source substrate that promotes each OpenIE
fact to a graph node and binds it to its source passage through
endpoint and source edges (\S\ref{sec:method:substrate}). At query
time, an entity-anchor extractor and a dense entry retriever
together induce a query-local subgraph $G_q$ over this substrate,
restricting propagation to passages and facts the question can
plausibly need. We then run a personalized PageRank on a derived
graph $H_q$ over local fact nodes; the resulting evidence flow
$\pi_q$ scores each passage by aggregating mass on its grounded
facts (\S\ref{sec:method:retrieval}). Because $\pi_q$ can still
underweight bridges several hops from any anchor, \methodname{}
mines the evidence the query still needs by following
relation-grounded edges to admit bridge and tail-entity passages,
and then applies a coverage-aware reranking step that selects a
final list under a noisy-OR coverage model over question-side
requirements (\S\ref{sec:method:enhancement}). The query-local
diffusion stage keeps retrieval focused, evidence-need mining
repairs missing bridges, and coverage-aware reranking prevents the
final list from collapsing into a single entity neighborhood. We
evaluate \methodname{} on HotpotQA, 2WikiMultiHopQA, and MuSiQue
under a unified GPT-4o-mini reader, comparing against classic text
RAG, GraphRAG, and evidence-enhanced GraphRAG baselines, and we
report retrieval quality, downstream QA, ablations of each
operator, a substrate isolation experiment, and an
upstream-extractor sensitivity check.

The main contributions are as follows.
\begin{itemize}
    \item We frame retrieval-first multi-hop QA as a conversion
    problem of turning graph-level relevance into a connected and
    covering evidence list under a small reader budget, identifying
    bridge underweighting and entity-neighborhood crowding as the
    two failure modes this conversion must address.
    \item We instantiate this framing in \methodname{} on top of a
    fact-source substrate that promotes each OpenIE fact to a graph
    node and binds it to its source passage. Query-local
    personalized PageRank, evidence-need mining, and noisy-OR
    coverage reranking jointly target focused propagation, bridge
    recovery, and requirement coverage.
    \item We show on HotpotQA, 2WikiMultiHopQA, and MuSiQue that
    this design improves retrieval and answer quality over classic
    text RAG, GraphRAG, and evidence-enhanced GraphRAG baselines
    under a unified reader, with component ablations, a substrate
    isolation experiment, and an upstream-extractor sensitivity
    check that preserves the ranking of the top systems.
\end{itemize}
